\section{Evaluation Methodology}
\label{sec:problem}
Our methodology follows the effect of an inference change from resource use
to the final answer. It separates four questions, summarized in
Table~\ref{tab:research-questions}. The implemented workflow is shown in
Figure~\ref{fig:research-system}.

\begin{table}[!htb]
\caption{Research questions and the evidence used to answer them.}
\label{tab:research-questions}
\centering
\begin{tabular}{@{}p{.7cm}p{3.6cm}p{2.7cm}@{}}
\toprule
 & Question & Evidence\\
\midrule
RQ1 & Does released memory let concurrent work progress? & Process and node memory, update progress\\[3pt]
RQ2 & Do backend settings and CPU allocations improve execution? & Standard backend, prefill and generation times\\[3pt]
RQ3 & Do individual answers and citations remain valid? & Paired QA checks, raw answers and citations\\[3pt]
RQ4 & Does a change improve the full episode, and can its timing be predicted? & Publication, request completion and prediction errors\\
\bottomrule
\end{tabular}
\end{table}

\begin{figure*}[t]
\centering
\begin{tikzpicture}[font=\small,>=Latex,
 methodbox/.style={draw=black!50,rounded corners=2pt,align=left,text width=4.25cm,minimum height=2.5cm,inner sep=7pt},
 flow/.style={->,line width=.8pt}]
\node[methodbox,fill=answerblue!9] (one) at (0,0) {\textbf{1. Isolated controls}\\
Fix model and CPU budget.\\
Run each job alone,\\then run both together.\\Measure phase times.};
\node[methodbox,fill=updatepurple!9] (two) at (5.45,0) {\textbf{2. Application episodes}\\
Set arrival times. Publish v2.\\
Retrieve the new evidence.\\Queue inference.\\Check answers and citations.};
\node[methodbox,fill=freshorange!12] (three) at (10.9,0) {\textbf{3. Prediction test}\\
Estimate phase times.\\Save predictions and hashes.\\
Test other documents\\and arrival times.};
\draw[flow] (one)--(two);\draw[flow] (two)--(three);
\node[draw=black!40,rounded corners,align=center,text width=15.2cm,minimum height=.9cm,fill=black!3] (impl) at (5.45,-2.0)
{\textbf{Implemented framework:} experiment controller, CPU affinity checks, versioned index worker,\\
request event logs, phase and resource samples, answer, citation and version checks, artifact collection};
\draw[flow] (one.south)--(one.south |- impl.north);
\draw[flow] (two.south)--(impl.north);
\draw[flow] (three.south)--(three.south |- impl.north);
\end{tikzpicture}
\caption{The evaluation workflow and the components we implemented. E1 applies
step 1, E2 applies step 2 and E3 applies step 3. Inputs and resource budgets
are fixed for each matched comparison. Timing and answer checks remain separate,
and prediction errors are retained without refitting.}
\label{fig:research-system}
\end{figure*}


\subsection{Evaluation Steps}
\textbf{1. Establish isolated controls.} Fix the model, backend, inputs and
resource budget. Run ingestion and inference alone, then together. Record
startup, memory use and phase times. A standard backend provides the execution
reference so that poor settings are not mistaken for an optimization opportunity.

\textbf{2. Measure the application episode.} Start an update while the current
request runs. Schedule a new question independently, wait for its required index
version, retrieve evidence and submit the request. Compare policies on matched
inputs and arrivals. Keep answer, citation and version checks separate from
latency. This step connects a local resource change to the user's request.

\textbf{3. Test the explanation.} Estimate phase durations from the earlier
runs, save the predictions and their hashes, then evaluate other document and
arrival combinations. Retain errors and failed answers without changing the
policy, labels or predictor to improve the result. E1, E2 and E3 instantiate
these steps with the experiments defined in Section~\ref{sec:budget}.

\subsection{Episode Measurements}
The current request keeps its original context while the update worker chunks
and embeds changed documents and publishes v2. A new question must retrieve
and cite the new evidence. Let $T_c$ be current completion, $T_p$ publication
and $T_f$ new-answer completion, measured from update arrival. We report all
three times and latency from the new question's own arrival. Finishing
generation does not mean that an answer passes the quality checks.

In the earlier sequential control, the new request starts after both current
completion and publication. Its timing is
\begin{equation}
T_f=\max(T_c,T_p)+T_{\mathrm{retrieve}}+T_{\mathrm{new\ request}}.
\label{eq:critical-path}
\end{equation}
For independent arrivals, retrieval waits for arrival and publication. The new
request then queues for the inference slot. The E3 estimate follows this order,
using measured phase durations. The order is a property of the experiment.
Whether the phase estimates predict other cases is tested separately.

We have no application SLA or independently reviewed utility score. We therefore
report timing and answer checks directly, without a weighted score chosen after
the results. Correct answers, supported citations, relevant versions and
abstention are assessed as separate outcomes.

\subsection{Implemented Components}
\label{sec:mechanisms}
We implemented an experiment controller that starts isolated and concurrent
jobs, schedules request arrivals independently and enforces the CPU allocation
for each process. It checks thread affinity during execution, saves timing
and resource samples, and collects outputs before removing its experiment pod.

The controller connects inference to a document worker that reuses unchanged
embeddings, publishes index versions atomically and retrieves evidence for the
new request. Separate checks record answer agreement, citation validity and
the retrieved version. The framework also saves predictions and their hashes
before the final timing test, along with the inputs, configuration and source
hashes needed to trace each result. These are the components we built around
llama.cpp, MiniLM and existing data and scoring tools. We do not claim that
these general capabilities were absent from all previous systems.

\paragraph{Earlier Distributed Precision Control}
A separate four-node pilot changes only stage 0 of Qwen2.5-14B between existing
Q4\_K\_M/Q3\_K\_M artifacts. At a token boundary it exports KV, destroys and
unmaps the old executor, loads the target, imports state and checks re-export
integrity. Replaying the full prefix through all stages checks target-model accuracy. This validates
a transition path rather than a new quantizer or selective repair algorithm.
The detailed lifecycle is retained in the separate artifact appendix.

\paragraph{CPU Execution and Fixed Attention Omission}
The later monolithic worker keeps the full Q4\_K\_M model loaded, distinguishes
readiness from request timing and clears KV before each request. Four threads
serve prefill and two serve decode, with an explicit prompt boundary so a
singleton prefill tail remains in the prefill phase. Each logical frame contains 64 tokens. The physical microbatch size controls
how the graph is split. Runtime repacking
changes the CPU weight layout for existing quantized kernels \citep{llamacpp}.
It neither changes the GGUF file nor introduces a new quantization method.

We test raw microbatches of 2, 8, 32 and 64, and repacked microbatches of 2 and 64. Larger batches and
layout changes are opt-in approximate execution configurations: the raw-mb2
worker reproduces reference logits at 64/565 tokens, but that equality is not
assumed for other configurations. Repacking can increase startup and anonymous
memory, so those costs are measured separately from warm requests.

The fixed mask omits attention sublayers 39--42, including their Q/K/V/O
operations and KV writes, while retaining every FFN. The mask is constant
for a session. Online reactivation is not implemented: missing KV and changed
downstream states prevent treating it as a free reversible action. Attention
dropping is established prior work \citep{attentiondrop}. Residual similarity suggests candidate layers to omit. Answer checks determine
whether a configuration passes.

\paragraph{What Each Change Tests}
Precision changes target memory use. Batching and repacking target execution time.
Attention omission removes work at potential quality cost. Update deferral
protects the current request by delaying freshness. Each is evaluated at its
own endpoint before being connected to utility. A successful implementation,
a smaller process counter and an acceptable application outcome are different
claims. We implement no controller that jointly selects these actions online.

